\documentclass[12pt,a4paper]{article}
\usepackage[utf8]{inputenc}
\usepackage[T1]{fontenc}
\usepackage{amsmath,amssymb,amsfonts}
\usepackage{mathtools}
\usepackage{booktabs}
\usepackage{array}
\usepackage{geometry}
\usepackage{hyperref}
\usepackage{cite}
\geometry{margin=2.5cm}

\newcommand{\MeV}{\,\text{MeV}}
\newcommand{\GeV}{\,\text{GeV}}
\newcommand{\eV}{\,\text{eV}}
\newcommand{\keV}{\,\text{keV}}
\newcommand{\fpi}{f_\pi}
\newcommand{\qqbar}{\langle\bar{q}q\rangle}
\newcommand{\zch}{\mathfrak{z}}

\title{Nuclear and hadronic predictions of the sBootstrap Lagrangian}

\author{A.~Rivero\thanks{Universidad de Zaragoza.
Manuscript prepared with AI assistance (Claude/Anthropic).
Continues the program of \cite{Rivero2005}.}}
\date{March 2026}

\begin{document}
\maketitle

\begin{abstract}
The sBootstrap $\mathcal{N}=1$ Lagrangian~\cite{paperIII}, constructed from
$\mathrm{SU}(3)$ SQCD at $N_f = N_c = 3$ with an O'Raifeartaigh
deformation, makes concrete predictions at the hadronic and nuclear scales.
We systematically derive calculational consequences that test the
framework against known nuclear physics.

The SUSY pion--muon relation $m_\pi^2 - m_\mu^2 = \fpi^2$ predicts a
SUSY-effective pion decay constant $\fpi^{\text{SUSY}} = 91.19\MeV$,
$1.0\%$ below the PDG value.  Using this value in the
Goldberger--Treiman relation reduces the discrepancy from
$1.3\%$ to $0.35\%$---a factor 3.8 improvement---and predicts
$g_{\pi NN} = 13.12$ ($0.3\%$ from experiment).  The quark condensate
derived from the SUSY Gell-Mann--Oakes--Renner relation gives
$|\qqbar|^{1/3} = 276\MeV$ ($6\%$ from the standard lattice value).

The meson spectrum has off-diagonal scalars at $\sqrt{2}\,\fpi \approx
130\MeV$ (pions) and diagonal scalars at $\fpi \approx 92\MeV$.
The diagonal mode, lighter than the pion, is identified with the
scalar-isoscalar channel and provides a longer-range nuclear attraction
($2.2\,\text{fm}$) than pion exchange ($1.4\,\text{fm}$), with
direct implications for the nuclear force at medium range.
Three mesino Dirac pairs at 11, 230, and 490$\eV$ predict a new
Yukawa force at $0.4$--$18\,\text{mm}$ range.  The muon mass is
derived as $m_\mu = \sqrt{m_\pi^2 - \fpi^2} = 104.9\MeV$ ($0.7\%$
from experiment), making it a prediction rather than a free parameter.
The Seiberg-dual sector predicts $m_d = 4.60\MeV$ ($0.1\sigma$)
and $m_s = 95.0\MeV$ ($1.9\sigma$).
\end{abstract}

%======================================================================
\section{Introduction}
%======================================================================

A supersymmetric Lagrangian at the nuclear scale was the implicit
promise of the programs that connected SUSY with hadronic
physics in the 1970s: from Miyazawa's baryon--meson
superalgebra~\cite{Miyazawa} through Catto and G\"ursey's
algebraic SUSY~\cite{CattoGursey} to the modern framework of
Seiberg duality~\cite{Seiberg1995} and ISS metastable
vacua~\cite{ISS2006}.  The sBootstrap~\cite{Rivero2005} makes this
promise concrete by providing an explicit $\mathcal{N}=1$
superpotential and K\"ahler potential from which the meson, baryon,
and mesino spectra follow.

This paper asks: given the Lagrangian of paper~[III]~\cite{paperIII},
what can be \emph{calculated} about nuclear and hadronic physics?
Not merely the mass spectrum (already derived in [III]), but the
nuclear force, the pion--nucleon coupling, the chiral condensate, the
Goldberger--Treiman relation, and other observables that have been
the standard tests of any theory of the strong interaction since the
1960s.

The central tool is the identification $\tilde{m}^2 = \fpi^2$: the
SUSY-breaking soft mass equals the pion decay constant squared.  This
single identification---fixing the scale of SUSY breaking at
$\fpi \approx 92\MeV$---turns the algebraic mass predictions of [III]
into dimensional predictions for hadronic observables.

Section~\ref{sec:scale} derives the SUSY pion decay constant from the
pion--muon relation.  Section~\ref{sec:mesons} computes the meson
spectrum and chiral condensate from the Seiberg vacuum.
Section~\ref{sec:nuclear} uses the resulting parameters to compute
the pion--nucleon coupling, the Goldberger--Treiman relation, and the
nuclear potential.  Section~\ref{sec:muon} discusses the muon as a
derived quantity and its implications for muonic atoms and
muon-catalyzed fusion.  Section~\ref{sec:supertrace} derives the
supertrace sum rule.  Section~\ref{sec:mesinos} describes the mesino
spectrum.  Section~\ref{sec:dualkoide} presents the Seiberg-dual
quark mass predictions.  Section~\ref{sec:summary} summarizes the
experimental programme.

%======================================================================
\section{The pion decay constant from SUSY}
\label{sec:scale}
%======================================================================

In the sBootstrap, the pion and the muon belong to the same SUSY
multiplet: the pion is a scalar meson $M^u_d$ and the muon is its
fermionic partner (a mesino in the $\mathrm{SU}(2)$ confining
sector). SUSY-breaking splits the masses:
\begin{equation}
  m_\pi^2 - m_\mu^2 = \fpi^2\,.
  \label{eq:pimu}
\end{equation}
This motivates defining a SUSY-effective pion decay constant:
\begin{equation}
  \fpi^{\text{SUSY}} \;\equiv\; \sqrt{m_\pi^2 - m_\mu^2}\,.
\end{equation}
Evaluating with PDG 2024~\cite{PDG2024}:
\begin{align}
  \fpi^{\text{SUSY}}
  &= \sqrt{(139.57)^2 - (105.66)^2}\MeV
   = \sqrt{8\,316}\MeV
   = 91.19\MeV\,,\\
  \fpi^{\text{PDG}} &= 92.07 \pm 0.57\MeV\,,
\end{align}
a $1.0\%$ discrepancy.  Conversely, assuming the relation is exact,
the muon mass is a \emph{derived} quantity:
\begin{equation}
  m_\mu^{\text{SUSY}} = \sqrt{m_\pi^2 - \fpi^2}
  = \sqrt{19\,480 - 8\,477}\MeV
  = 104.9\MeV\,,
  \label{eq:mu_pred}
\end{equation}
versus $m_\mu^{\text{exp}} = 105.66\MeV$ ($0.72\%$ discrepancy).
The muon mass is no longer a free parameter: it is determined by
$m_\pi$ and $\fpi$, both QCD observables.

The $\sim 1\%$ precision of~\eqref{eq:pimu} admits corrections from:
the electromagnetic pion mass (the DGMLY
contribution~\cite{DGMLY1967}
$m_{\pi^\pm}^2 - m_{\pi^0}^2 = 1{,}261\MeV^2$),
Kähler corrections to the off-diagonal meson mass, and
higher-loop CW contributions to the pseudo-modulus.

%======================================================================
\section{The meson spectrum and chiral condensate}
\label{sec:mesons}
%======================================================================

\subsection{Meson mass matrix}

At the Seiberg vacuum, the meson matrix $M$ acquires a diagonal VEV
$\langle M^i_j \rangle = v_i\delta^i_j$ with $v_i \propto 1/m_i$
(the seesaw).  The scalar meson fluctuations split into:

\begin{center}
\begin{tabular}{lccl}
\toprule
Sector & Count & Mass & QCD identification \\
\midrule
Off-diagonal $M^i_j$ & 12 real & $\sqrt{2}\,\fpi \approx 130\MeV$
  & $\pi^\pm$, $K^\pm$, $K^0$, $\bar K^0$ \\
Diagonal $\delta M^i_i$ & 6 real & $\fpi \approx 92\MeV$
  & $\pi^0$, $\eta$, $\eta'$, $\sigma$ \\
\bottomrule
\end{tabular}
\end{center}

The factor $\sqrt{2}$ between off-diagonal and diagonal masses
is the same ratio appearing in the O'Raifeartaigh $\sigma$--$\pi$
splitting at $t = \sqrt{3}$.

The pion mass prediction $m_\pi^{\text{pred}} = \sqrt{2}\,\fpi
= 130\MeV$ agrees with $m_{\pi^0} = 135\MeV$ to $3.7\%$
and with $m_{\pi^\pm} = 139.6\MeV$ to $6.9\%$.
The kaon masses ($494$--$498\MeV$) are off by a factor of $\sim 3.8$;
this reflects that the leading-order soft mass is flavor-universal
while the physical kaon mass is driven by the large $m_s/m_{u,d}$
ratio through the flavor-dependent Kähler potential.

\subsection{The diagonal mode as the scalar channel}

The six diagonal real modes at $\fpi \approx 92\MeV$ include the
scalar-isoscalar combination $(M^u_u + M^d_d)/\sqrt{2}$, which carries
the quantum numbers of the $\sigma$ meson ($f_0$, $I=0$, $J^P=0^+$).
At $92\MeV$, this mode is lighter than the pion:
\begin{equation}
  m_\sigma^{\text{sB}} = \fpi = 92\MeV\,,
  \qquad \text{vs.}\qquad
  m_\sigma^{\text{QCD}} = f_0(500)\text{: pole at }400\text{--}550\MeV\,.
\end{equation}
The sBootstrap sigma is lighter than the pion and much lighter than
the broad $f_0(500)$ resonance. Two interpretations are possible:
\begin{enumerate}
\item \textbf{Artifact}: the leading-order soft mass does not capture
the full dynamics of the scalar-isoscalar channel, which receives large
corrections from quark loops, $\eta$--$\eta'$ mixing, and the axial
anomaly.  In this case, $92\MeV$ is a tree-level
approximation that moves to $\sim 500\MeV$ after quantum corrections.
\item \textbf{Prediction}: there exists a narrow scalar state near
$92\MeV$ that has escaped detection because its width is suppressed
by SUSY relations.  This would be a dramatic experimental prediction.
\end{enumerate}
We consider interpretation~(1) more likely, since the scalar-isoscalar
channel is notorious for large non-perturbative corrections.  However,
the Yukawa range $\hbar c / m_\sigma^{\text{sB}} = 2.16\,\text{fm}$
is the correct order of magnitude for the medium-range nuclear
attraction (Section~\ref{sec:nuclear}).

\subsection{The GOR relation and the chiral condensate}

The Gell-Mann--Oakes--Renner relation~\cite{GOR1968}
\begin{equation}
  m_\pi^2 \fpi^2 = -(m_u + m_d)\,\qqbar
\end{equation}
connects the pion mass to the quark condensate.
In the sBootstrap, the off-diagonal meson mass-squared is $2\fpi^2$
(the soft mass). Substituting:
\begin{equation}
  -(m_u + m_d)\,\qqbar = 2\fpi^4\,,
  \label{eq:qqbar_susy}
\end{equation}
which gives a parameter-free prediction for the chiral condensate:
\begin{equation}
  |\qqbar|^{1/3}
  = \left(\frac{2\fpi^4}{m_u + m_d}\right)^{1/3}
  = \left(\frac{2\times(92.07)^4}{6.83}\right)^{1/3}\MeV
  = 276\MeV\,.
\end{equation}
The standard lattice value is $|\qqbar|^{1/3} \approx 260\MeV$
(at $\mu = 2\GeV$, $\overline{\text{MS}}$), a $6.2\%$ discrepancy.
This is consistent given the leading-order approximation: the
$\sqrt{2}$ factor in $m_\pi^{\text{sB}} = \sqrt{2}\fpi$ versus the
physical $m_\pi = 135\MeV$ propagates into the condensate estimate.
With the physical pion mass, GOR gives
$|\qqbar|^{1/3} = (m_\pi^2 \fpi^2 / (m_u+m_d))^{1/3} = 261\MeV$,
recovering the standard value.

%======================================================================
\section{The nuclear force: GT, $g_{\pi NN}$, and the NN potential}
\label{sec:nuclear}
%======================================================================

\subsection{The Goldberger--Treiman relation}

The Goldberger--Treiman (GT) relation~\cite{GT1958}
\begin{equation}
  g_{\pi NN}\,\fpi = g_A\,m_N
  \label{eq:GT}
\end{equation}
connects the pion--nucleon coupling $g_{\pi NN}$ to the axial coupling
$g_A = 1.2756 \pm 0.0013$ (neutron $\beta$ decay) and the nucleon mass
$m_N = 938.3\MeV$.  With $g_{\pi NN} = 13.17 \pm 0.06$ (Nijmegen
partial-wave analysis), the GT discrepancy is
\begin{equation}
  \Delta_{\text{GT}} \equiv 1 - \frac{g_A m_N}{g_{\pi NN}\fpi}\,.
\end{equation}

Using the PDG pion decay constant $\fpi = 92.07\MeV$:
\begin{equation}
  \Delta_{\text{GT}}^{\text{PDG}}
  = 1 - \frac{1.2756 \times 938.3}{13.17 \times 92.07}
  = 1 - 0.987 = 1.3\%\,.
\end{equation}
Using the SUSY-predicted value $\fpi^{\text{SUSY}} = 91.19\MeV$:
\begin{equation}
  \boxed{
  \Delta_{\text{GT}}^{\text{SUSY}}
  = 1 - \frac{1.2756 \times 938.3}{13.17 \times 91.19}
  = 1 - 0.997 = 0.35\%\,.
  }
\end{equation}
The SUSY-effective pion decay constant reduces the GT discrepancy by a
factor of~$3.8$.  This is the single most concrete nuclear prediction
of the sBootstrap: the ``correct'' $\fpi$ for the GT relation is
$\sqrt{m_\pi^2 - m_\mu^2}$, not the PDG value from $\pi \to \mu\nu$.

\subsection{The pion--nucleon coupling}

Alternatively, using $g_A$ and $m_N$ as inputs and the GT relation as
an \emph{exact} consequence of SUSY, we predict $g_{\pi NN}$:
\begin{align}
  g_{\pi NN}^{\text{PDG}}
  &= \frac{g_A m_N}{\fpi^{\text{PDG}}}
   = \frac{1.2756 \times 938.3}{92.07} = 13.00\,,\\[4pt]
  g_{\pi NN}^{\text{SUSY}}
  &= \frac{g_A m_N}{\fpi^{\text{SUSY}}}
   = \frac{1.2756 \times 938.3}{91.19} = 13.12\,.
\end{align}
The experimental value $g_{\pi NN} = 13.17 \pm 0.06$ is reproduced
to $0.38\%$ by $\fpi^{\text{SUSY}}$, versus $1.3\%$ with
$\fpi^{\text{PDG}}$.

\subsection{The one-pion-exchange nuclear potential}

The backbone of nuclear physics is the one-pion-exchange potential
(OPEP).  For the $NN$ central channel at distance $r$:
\begin{equation}
  V_{\text{OPEP}}^{\text{(central)}}(r)
  = -\frac{g_{\pi NN}^2}{4\pi}\,
    \frac{m_\pi^3}{12\,m_N^2}\,
    \frac{e^{-m_\pi r / (\hbar c)}}{m_\pi r / (\hbar c)}\,.
\end{equation}
At $r = 1\,\text{fm}$, with $g_{\pi NN} = 13.12$ (SUSY) and
$m_\pi = 139.6\MeV$:
\begin{equation}
  V_{\text{OPEP}}(1\,\text{fm}) = -2.5\MeV\,.
\end{equation}
The pion exchange range is $\hbar c / m_\pi = 1.41\,\text{fm}$.

\subsection{Medium-range attraction: the diagonal meson exchange}

The sBootstrap predicts a scalar-isoscalar mode at $m_\sigma^{\text{sB}}
= \fpi = 92\MeV$.  If this mode couples to nucleons with strength
$g_{\sigma NN}$, it contributes a central attractive potential
\begin{equation}
  V_\sigma(r) = -\frac{g_{\sigma NN}^2}{4\pi}\,
    \frac{e^{-m_\sigma r/(\hbar c)}}{r}\,,
\end{equation}
with range $\hbar c / m_\sigma^{\text{sB}} = 2.16\,\text{fm}$---longer
than the pion range by a factor $1.53$.

In the standard nuclear force paradigm, the medium-range attraction
(the ``sigma channel'') at $r \approx 1$--$2\,\text{fm}$ is
provided by two-pion exchange or an effective $\sigma$ meson at
$400$--$550\MeV$ (range $0.35$--$0.49\,\text{fm}$).  The sBootstrap
places this attraction at a larger range.  The Yukawa factor at
$r = 1\,\text{fm}$ is $e^{-x}/x = 1.36$ for $m_\sigma = 92\MeV$,
versus $0.045$ for $m_\sigma = 450\MeV$: a factor $30\times$ stronger.
Even a small coupling $g_{\sigma NN} \sim 1$ would produce a significant
medium-range potential.

Whether this longer-range scalar exchange improves or worsens the
nuclear force fit is an open question.  It may be that the leading-order
$92\MeV$ mode receives large corrections (as argued in
Section~\ref{sec:mesons}) and moves to $\sim 500\MeV$ in the full
theory.  Alternatively, the extra attraction at $2\,\text{fm}$ could
account for the intermediate-range $NN$ attraction that in the
Bonn/CD-Bonn potentials is parameterized by correlated two-pion
exchange with an effective mass of $\sim 300$--$500\MeV$; the sBootstrap
diagonal mode provides a candidate for the physical origin of this term
at a lower mass.

\subsection{The deuteron binding: qualitative estimate}

The deuteron binding energy $E_d = 2.224\MeV$ corresponds to a
characteristic size $1/\kappa = 6.1\,\text{fm}$
($\kappa = \sqrt{m_N E_d/2}/\hbar c = 0.164\,\text{fm}^{-1}$).
This is $4.3$ times the pion range and $2.8$ times the sBootstrap
sigma range: the deuteron ``lives'' in the tail of the nuclear potential,
dominated by the lightest exchanged meson.

If the lightest scalar exchange has mass $m_\sigma^{\text{sB}} = 92\MeV$,
the Yukawa tail at the deuteron radius is $e^{-0.92 \cdot 6.1/2.16}
= e^{-2.6} = 0.074$---not negligible.  At the pion mass,
$e^{-6.1/1.41} = e^{-4.3} = 0.013$---exponentially suppressed.
This suggests that the deuteron binding in the sBootstrap is sensitive
to the scalar-isoscalar channel at $92\MeV$, and getting $E_d$ right
constrains $g_{\sigma NN}$.

A full calculation of $E_d$ from the sBootstrap requires solving the
Schr\"odinger equation with the tensor and central OPEP plus the
scalar exchange, and lies beyond the scope of this paper.  However,
the qualitative point is that the sBootstrap provides \emph{both}
the pion and sigma parameters from a single Lagrangian---precisely the
information needed to compute nuclear binding.

\subsection{The $dt$ astrophysical $S$-factor}

The nuclear cross-section at low energy is parameterized by the
astrophysical $S$-factor:
\begin{equation}
  \sigma(E) = \frac{S(E)}{E}\,\exp\!\left(-\sqrt{E_G/E}\right),
\end{equation}
where $E_G = 2\pi^2 \alpha^2 \mu_N c^2 \approx 0.50\MeV$ is the
Gamow energy for $d$--$t$ ($\mu_N$ the nuclear reduced mass,
$Z_1 = Z_2 = 1$).  For the reaction
$d(t,n){}^4\text{He}$, the $S$-factor at low energy is dominated
by the $J^\pi = 3/2^+$ resonance in ${}^5\text{He}$ at
$E_R \approx 64\keV$ above the $d$--$t$ threshold.

The $S$-factor is well measured above $\sim 5\keV$, but stellar
burning in the solar core ($T \approx 1.5\keV$) and the margins of
inertial confinement fusion (ICF) ignition ($T \sim 5$--$10\keV$)
rely on theoretical extrapolation.  The R-matrix analysis
parameterizes $S(E)$ in terms of resonance energies and reduced
widths, which in turn depend on the nuclear potential at the
channel radius.

The sBootstrap modifies two potential ingredients:
\begin{itemize}
\item OPEP with $g_{\pi NN} = 13.12$ (versus the standard
  $13.17$): a $0.76\%$ change in $g_{\pi NN}^2$, hence a $0.76\%$
  shift in the tensor and central OPEP.
\item The scalar-isoscalar exchange at $92\MeV$ (range $2.16\,\text{fm}$)
  instead of the phenomenological $\sigma$ at $\sim 500\MeV$
  (range $0.4\,\text{fm}$).
\end{itemize}
The OPEP correction is small.  The sigma modification is potentially
large: the $dt$ scattering wave function at the channel radius
($r_c \approx 5\,\text{fm}$) is exponentially sensitive to the
potential tail.  A Yukawa tail $e^{-r/2.16\,\text{fm}}$ at
$r = 5\,\text{fm}$ gives $e^{-2.3} = 0.10$, whereas
$e^{-r/0.4\,\text{fm}} = e^{-12.5} \approx 4 \times 10^{-6}$:
the sBootstrap sigma is a factor $\sim 2.5 \times 10^4$ stronger
at the channel radius.

A quantitative $S$-factor calculation requires an R-matrix analysis
with the sBootstrap potential.  The prediction is that the low-energy
$S(E)$ extrapolation differs from the standard potential-model result
due to the longer-range scalar attraction.  At stellar energies
($E \sim 1\keV$), where the Gamow penetration factor
$e^{-\sqrt{E_G/E}} \sim e^{-22}$ dominates, even a modest change
in the pre-exponential factor is significant for stellar
nucleosynthesis rates.

\subsection{The solar $pp$ $S$-factor}

The $pp$ reaction $p + p \to d + e^+ + \nu_e$ proceeds through the
weak interaction and has an astrophysical $S$-factor
$S_{11}(0) = (4.01 \pm 0.03) \times 10^{-25}\MeV\cdot\text{barn}$~\cite{AdelbergerSFactor2011}.
The $0.7\%$ theoretical uncertainty is dominated by the
overlap of the $pp$ scattering wave function with the deuteron
bound-state wave function at zero separation.

Both wave functions are determined by the $NN$ potential.  The $pp$
scattering wave function at short range depends on the strong
interaction, specifically on the scattering length $a_{pp}$ and
effective range $r_{pp}$.  The deuteron wave function's asymptotic
normalization depends on the potential tail.

The sBootstrap contribution is twofold:
\begin{enumerate}
\item The potential parameters ($g_{\pi NN}$, $m_\pi$, $m_\sigma$,
  $g_{\sigma NN}$) are fixed rather than fitted.  This eliminates
  the model dependence of the nuclear wave functions.
\item The longer-range sigma exchange modifies the wave function
  overlap integral.  The effect is largest at the nuclear surface
  ($r \sim 2\,\text{fm}$), where the sigma Yukawa tail is significant.
\end{enumerate}

A $0.35\%$ shift in $g_{\pi NN}$ propagates into $\sim 0.7\%$ in
the OPEP potential, comparable to the current theoretical uncertainty
in $S_{11}$.  If the sBootstrap potential gives a different wave
function overlap, the solar $pp$ rate---and hence the predicted
solar neutrino flux---shifts correspondingly.  The solar neutrino
measurements by SNO, Borexino, and Super-Kamiokande currently
constrain $S_{11}$ to $\sim 1\%$; a competitive $S_{11}$ from the
sBootstrap potential would be a stringent test.

\subsection{The $\alpha$--$\alpha$ potential and the Hoyle state}
\label{subsec:hoyle}

The triple-alpha process $3\alpha \to {}^{12}\text{C}$ is the
bottleneck of stellar helium burning and the origin of all
terrestrial carbon.  It proceeds through a two-step resonance:
\begin{equation}
  \alpha + \alpha \to {}^8\text{Be}^*\,,\qquad
  {}^8\text{Be}^* + \alpha \to {}^{12}\text{C}^*
  \;(\text{Hoyle state, }7.654\MeV)\,.
\end{equation}
The Hoyle state lies only $\Delta E = 379.47 \pm 0.15\keV$ above the
$3\alpha$ threshold.  This near-degeneracy has long been considered
a fine-tuning: a $\sim 4\%$ shift in $\Delta E$ would suppress
carbon production by orders of magnitude, preventing carbon-based
life~\cite{Hoyle1954,OberhumerCSoto2000}.

The Hoyle state energy depends on the nuclear potential at the
$\alpha$-cluster scale ($r \sim 3$--$4\,\text{fm}$): the
$\alpha$--$\alpha$ potential at medium range determines the energy
of the $0^+$ resonance in ${}^{12}\text{C}$.

The sBootstrap enters through the scalar-isoscalar exchange.  The
standard nuclear force models use a phenomenological $\sigma$ at
$400$--$550\MeV$ (range $0.35$--$0.5\,\text{fm}$), which is
essentially a contact interaction at the $\alpha$-cluster scale.
The sBootstrap diagonal mode at $92\MeV$ has range $2.16\,\text{fm}$
--- comparable to the $\alpha$--$\alpha$ separation in ${}^{12}\text{C}$
($\sim 3\,\text{fm}$).  At this distance:
\begin{align}
  \text{sBootstrap:}\quad &
    \frac{e^{-3.0/2.16}}{3.0/2.16} = \frac{e^{-1.39}}{1.39}
    = \frac{0.249}{1.39} = 0.179\,,\\
  \text{Standard:}\quad &
    \frac{e^{-3.0/0.44}}{3.0/0.44} = \frac{e^{-6.8}}{6.8}
    = \frac{0.0011}{6.8} = 1.6 \times 10^{-4}\,.
\end{align}
The sBootstrap sigma Yukawa factor at the $\alpha$-cluster distance
is $\sim 1{,}100$ times larger than the standard sigma.

This means the sBootstrap scalar exchange provides a qualitatively
different medium-range $\alpha$--$\alpha$ attraction: not a
short-range contact term but a finite-range Yukawa with substantial
strength at $3\,\text{fm}$.  The Hoyle state energy is determined
by the depth and shape of the $\alpha$--$\alpha$ potential well in
this regime.

The provocative possibility is that the Hoyle state near-degeneracy
($\Delta E / E_{\text{Hoyle}} = 5.0\%$) is not a fine-tuning but
a \emph{consequence} of the sBootstrap scalar mass: the $92\MeV$
diagonal mode, through its long-range $\alpha$--$\alpha$ attraction,
naturally places the Hoyle state close to threshold.  Testing this
requires an $\alpha$-cluster model calculation with the sBootstrap
potential, which we leave to future work.  But the numerology is
suggestive: the Hoyle state sits at the $\alpha$-cluster distance
($\sim 3\,\text{fm}$) set by the sBootstrap sigma range
($2.16\,\text{fm}$), not by the standard sigma range
($0.4\,\text{fm}$).

\subsection{Nuclear equation of state at high compression}

In ICF, DT fuel is compressed to
$\rho \sim 300$--$1000\,\text{g/cm}^3$ (corresponding to
$100$--$350\times$ liquid density) before ignition.  At these
densities, the average internuclear separation is
\begin{equation}
  d \sim \left(\frac{m_N}{\rho}\right)^{1/3}
  = \left(\frac{1.67 \times 10^{-24}\,\text{g}}{300\,\text{g/cm}^3}\right)^{1/3}
  \approx 1.8\,\text{fm}\,,
\end{equation}
comparable to the sBootstrap sigma range ($2.16\,\text{fm}$).
The nuclear equation of state (EOS) at these densities depends on
the $NN$ potential: repulsion at short range ($< 0.5\,\text{fm}$)
stiffens the EOS, while attraction at medium range softens it.

The standard nuclear EOS models (Skyrme, relativistic mean field)
parameterize the medium-range attraction with coupling constants
fitted to nuclear binding energies and saturation properties.  The
sBootstrap fixes these parameters from the Lagrangian:
\begin{itemize}
\item Saturation density: determined by the balance of OPEP
  (tensor repulsion at short range) and sigma attraction.  A lighter
  sigma pushes saturation to lower density (the attraction turns on
  at larger $r$).
\item Compression modulus: $K \approx 240 \pm 20\MeV$.  The
  sBootstrap sigma at $92\MeV$ gives a softer EOS near
  saturation, potentially reducing $K$.
\item Symmetry energy: dominated by the isovector channel ($\rho$
  meson exchange), which the sBootstrap does not modify at leading
  order.
\end{itemize}

For ICF, a softer EOS means less pressure is needed to reach a given
compression, but the ignition threshold also shifts.  The net effect
on ICF performance depends on the interplay between compression and
heating, which requires radiation-hydrodynamics simulation with the
sBootstrap EOS.  The prediction is that the sBootstrap EOS differs
from standard Skyrme parameterizations at $\rho \gtrsim
100\,\text{g/cm}^3$, where the internuclear separation enters the
sigma exchange range.

%======================================================================
\section{The muon mass and muon-catalyzed fusion}
\label{sec:muon}
%======================================================================

Muon-catalyzed fusion (MCF) is the process by which a negative muon,
substituting for an electron in a hydrogen isotope atom, brings two
nuclei close enough to fuse at room temperature.  The MCF cycle
depends on a chain of atomic and nuclear parameters---the muonic Bohr
radius, the molecular ion binding, the nuclear tunneling probability,
and the sticking fraction---all of which, in the sBootstrap, trace
back to two QCD observables: $m_\pi$ and $\fpi$.

\subsection{The muon as a derived quantity}

In the sBootstrap, the muon mass is not free: from~\eqref{eq:mu_pred},
\begin{equation}
  m_\mu^{\text{SUSY}} = \sqrt{m_\pi^2 - \fpi^2} = 104.9\MeV\,,
\end{equation}
a $0.72\%$ prediction ($0.76\MeV$ below the experimental value
$105.66\MeV$).  This means the mass hierarchy that makes MCF
possible is a \emph{consequence} of the SUSY-breaking structure:
\begin{equation}
  \frac{m_\mu}{m_\pi} = \sqrt{1 - \frac{\fpi^2}{m_\pi^2}}
  = 0.752\,,
\end{equation}
using $\fpi^2/m_\pi^2 \approx 0.435$.
A heavier muon ($m_\mu > m_\pi$) would make MCF
impossible; the sBootstrap predicts $m_\mu < m_\pi$
as a structural necessity of SUSY breaking.

\subsection{The MCF cycle: step by step}

The $dt$ fusion cycle proceeds through four stages.  At each stage
we identify the sBootstrap parameter that controls the physics.

\paragraph{Step 1: Muonic atom formation.}
A negative muon is captured by a deuteron, forming muonic deuterium
$d\mu$ with Bohr radius
\begin{equation}
  a_{d\mu} = \frac{\hbar c}{\alpha\,m_{\text{red}}}
  = \frac{197.3\MeV\cdot\text{fm}}{(1/137)\times 100.0\MeV}
  = 270.3\,\text{fm}\,,
\end{equation}
where $m_{\text{red}} = m_\mu m_d/(m_\mu + m_d) = 100.0\MeV$ is the
$\mu$--$d$ reduced mass.  This is $196$ times smaller than the
electronic Bohr radius ($53{,}000\,\text{fm}$).  The 1s binding energy
is $E_{1s} = m_{\text{red}}\alpha^2/2 = 2.66\keV$.

With $m_\mu^{\text{SUSY}}$: $m_{\text{red}}^{\text{SUSY}} = 99.3\MeV$,
$a_{d\mu}^{\text{SUSY}} = 272.2\,\text{fm}$ (shift $+1.9\,\text{fm}$,
$+0.69\%$), and $E_{1s}^{\text{SUSY}} = 2.64\keV$ (shift $-18\eV$).

\emph{sBootstrap parameter}: $m_\mu = \sqrt{m_\pi^2 - \fpi^2}$ sets
the atomic scale.

\paragraph{Step 2: Molecular ion formation (Vesman resonance).}
The $d\mu$ atom collides with a $\text{DT}$ molecule and resonantly
forms the $dt\mu$ molecular ion in the loosely bound
$(J=1,\,v=1)$ state, with binding energy
$\varepsilon_{11} = 0.660\eV$ below the $d\mu(1s) + t$ threshold.
The Vesman mechanism~\cite{Vesman1967} requires near-coincidence
between $\varepsilon_{11}$ and the rovibrational excitation energy of
the host $\text{DT}$ molecule: the released binding energy excites
a rovibrational quantum, and the resonance condition is
\begin{equation}
  E_{\text{kin}} + \varepsilon_{11}
  = \Delta E_{\text{rovib}}(\text{DT})\,.
\end{equation}
At $T \approx 300\,\text{K}$, the thermal energy $kT \approx 26\,\text{meV}$
provides the kinetic energy for resonance matching.  The formation
rate reaches $\lambda_{dt\mu} \sim 10^{8}\,\text{s}^{-1}$ at
optimal temperature---the rate-limiting step of the MCF cycle.

The binding energy $\varepsilon_{11}$ scales with the $\mu$--nucleus
reduced mass.  With $m_\mu^{\text{SUSY}}$:
\begin{equation}
  \varepsilon_{11}^{\text{SUSY}}
  = \varepsilon_{11} \times
    \frac{m_{\text{red}}^{\text{SUSY}}}{m_{\text{red}}^{\text{exp}}}
  = 0.660 \times \frac{99.3}{100.0}\eV
  = 0.655\eV\,.
\end{equation}
The shift $\Delta\varepsilon_{11} = -4.5\,\text{meV}$ is $17\%$
of the thermal energy at $300\,\text{K}$.  This does not destroy
the resonance (which has thermal width $\sim kT$) but shifts
the optimal temperature: to re-centre the resonance peak,
the temperature would need to increase by roughly
\begin{equation}
  \Delta T \approx
  \frac{\Delta\varepsilon_{11}}{k} \approx 52\,\text{K}\,,
\end{equation}
from $\sim 300\,\text{K}$ to $\sim 350\,\text{K}$.

\emph{sBootstrap parameter}: $\varepsilon_{11}$ through $m_\mu$;
the Vesman resonance condition is a derived consequence.

\paragraph{Step 3: Nuclear fusion.}
Inside the $dt\mu$ molecular ion, the $d$ and $t$ nuclei are held at
an equilibrium distance $R_{\text{eq}} \sim 0.6\,a_{d\mu}
\approx 160\,\text{fm}$---close enough for the nuclear wave functions
to overlap through Coulomb tunnelling.  The fusion reaction
\begin{equation}
  d + t \;\to\; {}^4\text{He}\;(3.54\MeV)
  \;+\; n\;(14.05\MeV)
\end{equation}
releases $Q_{dt} = 17.59\MeV$.  The fusion rate inside
the molecular ion is $\lambda_{\text{fus}} \sim 10^{12}\,\text{s}^{-1}$,
essentially instantaneous compared to the molecular formation rate.
Fusion is not the bottleneck.

The fusion rate depends on the nuclear $S$-factor $S(E)$ for
$d(t,n)\alpha$ at low energy, which in turn depends on the nuclear
potential.  At internuclear distances $r > 1\,\text{fm}$, the
potential is dominated by one-pion exchange (Section~\ref{sec:nuclear})
with $g_{\pi NN}$.  The sBootstrap value $g_{\pi NN}^{\text{SUSY}}
= 13.12$ (from the GT relation) modifies the OPEP by $(13.12/13.17)^2
= 0.992$, a $0.8\%$ change in the potential.  Since
$\lambda_{\text{fus}} \gg \lambda_{dt\mu}$, this does not limit the
MCF cycle but is relevant for the nuclear $S$-factor itself.

\emph{sBootstrap parameter}: $g_{\pi NN}$ from
$g_A m_N / \fpi^{\text{SUSY}}$ enters the tunnelling matrix element.

\paragraph{Step 4: Muon release and sticking.}
After fusion, the $\alpha$ particle recoils at
$v_\alpha/c = 0.044$ ($E_\alpha = 3.54\MeV$).
The muon, initially in the molecular orbital, can either:
\begin{itemize}
\item[(a)] escape and catalyse another fusion (probability $1-\omega_s$), or
\item[(b)] stick to the $\alpha$ particle, forming $\alpha\mu^+$
  (probability $\omega_s$).
\end{itemize}
The initial sticking probability is $\omega_s^0 \approx 0.9\%$;
after reactivation (stripping of $\alpha\mu^+$ in subsequent collisions),
the effective sticking is $\omega_s^{\text{eff}} \approx 0.5$--$0.7\%$.

The sticking probability in the sudden approximation depends on the
overlap integral
\begin{equation}
  \omega_s \propto
  \bigl|\langle\psi_{\alpha\mu}|\,e^{i\mathbf{p}_\mu\cdot\mathbf{r}/\hbar}\,
  |\psi_{\text{mol}}\rangle\bigr|^2\,,
\end{equation}
where $\mathbf{p}_\mu = m_\mu \mathbf{v}_\alpha$ is the momentum
kick imparted to the muon by the recoiling $\alpha$.  The
adiabaticity parameter
$\xi = p_\mu a_{\alpha\mu} / \hbar
= m_\mu v_\alpha a_{\alpha\mu} / \hbar c
\approx 3.15$ places the system in the intermediate regime
(neither sudden nor adiabatic).

The sticking probability scales approximately as $m_{\text{red}}^3$
(from the three-dimensional momentum-space overlap).  With
$m_\mu^{\text{SUSY}}$:
\begin{equation}
  \omega_s^{\text{SUSY}}
  = \omega_s^{\text{exp}} \times
    \left(\frac{m_{\text{red}}^{\text{SUSY}}}{m_{\text{red}}^{\text{exp}}}\right)^3
  = 0.67\% \times \left(\frac{99.3}{100.0}\right)^3
  = 0.656\%\,,
\end{equation}
a $2.0\%$ reduction in the sticking probability.

\emph{sBootstrap parameter}: $\omega_s$ through $m_\mu$ and the
momentum kick $p_\mu = m_\mu v_\alpha$.

\subsection{Fusions per muon and the break-even problem}

The number of fusions catalysed by a single muon before it decays
($\tau_\mu = 2.197\,\mu\text{s}$) or sticks is
\begin{equation}
  N_{\text{fus}} = \frac{\lambda_c\,\tau_\mu}
  {1 + \omega_s\,\lambda_c\,\tau_\mu}\,,
\end{equation}
where $\lambda_c \sim 10^8\,\text{s}^{-1}$ is the cycling rate
(dominated by the molecular formation step).

In the sticking-limited regime ($\omega_s\lambda_c\tau_\mu \gg 1$):
$N_{\text{fus}} \approx 1/\omega_s$.  With the SUSY-corrected
sticking:

\begin{center}
\begin{tabular}{lcc}
\toprule
 & Experimental $m_\mu$ & SUSY $m_\mu$ \\
\midrule
$\omega_s^{\text{eff}}$ & $0.670\%$ & $0.656\%$ \\
$N_{\text{fus}}$ (sticking limit) & 149 & 152 \\
$N_{\text{fus}}$ (full, $\lambda_c = 10^8\,\text{s}^{-1}$) & 89 & 90 \\
\bottomrule
\end{tabular}
\end{center}

The gain from the SUSY correction is $+3$ fusions per muon in the
sticking limit---a modest $2\%$ improvement.  Break-even for energy
production requires $N_{\text{fus}} \geq E_{\mu\text{-prod}}/Q_{dt}
\approx 5000/17.6 \approx 284$ fusions per muon, assuming
$5\GeV$ to produce one muon.  The sBootstrap correction closes
about $2\%$ of the gap.

The significance is not in the magnitude of the correction but in
its \emph{calculability}: in the sBootstrap, every parameter entering
the MCF cycle---$m_\mu$, $a_{d\mu}$, $\varepsilon_{11}$,
$g_{\pi NN}$, $\omega_s$---is determined by $(m_\pi, \fpi, g_A, m_N)$.
The entire MCF cycle becomes a derived consequence of the SUSY
Lagrangian plus QCD inputs.

\subsection{The complete SUSY chain for MCF}

The sBootstrap provides a single causal chain from the QCD Lagrangian
to the MCF fusion count:
\begin{equation}
\begin{array}{rcll}
  m_\pi,\,\fpi & \xrightarrow{\text{SUSY}} & m_\mu = \sqrt{m_\pi^2 - \fpi^2}
    & \text{(muon mass)} \\
  m_\mu & \xrightarrow{\text{QED}} & a_{d\mu} = \hbar c/(\alpha\,m_{\text{red}})
    = 270\,\text{fm} & \text{(muonic atom)} \\
  a_{d\mu} & \to & \varepsilon_{11} = 0.66\eV
    & \text{(molecular ion)} \\
  \varepsilon_{11} + kT & \to & \lambda_{dt\mu} \sim 10^8\,\text{s}^{-1}
    & \text{(Vesman resonance)} \\
  g_A,\,m_N,\,\fpi & \xrightarrow{\text{GT}} & g_{\pi NN} = 13.12
    & \text{(nuclear coupling)} \\
  g_{\pi NN},\,m_\pi & \to & V_{\text{OPEP}}(r)
    & \text{(nuclear potential)} \\
  V_{\text{OPEP}} & \to & \lambda_{\text{fus}} \sim 10^{12}\,\text{s}^{-1}
    & \text{(fusion rate)} \\
  m_\mu,\,v_\alpha & \to & \omega_s = 0.66\%
    & \text{(sticking fraction)} \\
  \lambda_c,\,\omega_s,\,\tau_\mu & \to & N_{\text{fus}} \approx 150
    & \text{(fusions/muon)}
\end{array}
\end{equation}
No step in this chain requires a free parameter beyond those fixed by
the sBootstrap Lagrangian and measured QCD/electroweak inputs.

\subsection{The pion--muon mass difference and pion lifetime}

The mass difference $m_\pi - m_\mu = 33.9\MeV$ is the available
energy for pion decay $\pi^+ \to \mu^+ \nu_\mu$; its largeness
(relative to $m_e$) determines the pion lifetime and hence the
range of the nuclear force in matter.  In the sBootstrap:
\begin{equation}
  m_\pi - m_\mu
  = m_\pi\left(1 - \sqrt{1 - \frac{\fpi^2}{m_\pi^2}}\right)
  \approx \frac{\fpi^2}{2\,m_\pi}
  = \frac{(92.1)^2}{2 \times 139.6}\MeV
  = 30.4\MeV\,,
\end{equation}
versus the observed $33.9\MeV$ ($10\%$ discrepancy, from
the $1\%$ error in $\fpi^{\text{SUSY}}/\fpi^{\text{PDG}}$ amplified
by the nonlinear $\sqrt{1-x}$ function).  The pion decay phase space
$\propto (m_\pi^2 - m_\mu^2)^2 = \fpi^4$ in the SUSY limit:
the pion lifetime is set by the fourth power of the SUSY-breaking
scale.

The pion lifetime controls the temporal reach of the nuclear force
in nuclear reactions and astrophysical processes: a shorter-lived
pion means a shorter effective range in time-dependent processes.
The sBootstrap connects this lifetime to $\fpi$ through the
pion--muon superpartnership.

%======================================================================
\section{The supertrace sum rule}
\label{sec:supertrace}
%======================================================================

The supertrace of the squared mass matrix at the Seiberg vacuum
receives contributions only from the soft SUSY-breaking masses
(the $W_{IJ}^2$ piece cancels between scalar and fermion sectors):
\begin{equation}
  \boxed{
    \text{STr}[M^2] = 18\,\fpi^2 = 152{,}352\MeV^2\,.
  }
\end{equation}
The factor $18 = 2 \times 9$: two real components per complex meson,
nine diagonal entries of the $3 \times 3$ meson matrix.  This is an
exact, parameter-free prediction.

In terms of hadronic spectroscopy:
\begin{equation}
  \sum_{\text{scalar mesons}} m_S^2
  - 2\sum_{\text{mesinos}} m_F^2
  = 18\,\fpi^2\,.
\end{equation}
With 12 off-diagonal scalars at $2\fpi^2$ each and 6 diagonal scalars
at $\fpi^2$ each:
$\sum m_S^2 = 12 \times 2\fpi^2 + 6 \times \fpi^2 = 30\,\fpi^2$.
The mesino contribution (Section~\ref{sec:mesinos}) is
$\sum m_F^2 \approx 0.3\MeV^2$, negligible.  So
$30\fpi^2 - 12\fpi^2 = 18\fpi^2$ after subtracting the fermion
mass contribution from the supertrace.

Testing this requires measuring the complete low-lying meson scalar
spectrum and comparing with $18\fpi^2$.  The pion--muon relation
already tests the off-diagonal piece; the diagonal piece at
$\fpi \approx 92\MeV$ is the new content.

%======================================================================
\section{The mesino spectrum}
\label{sec:mesinos}
%======================================================================

\subsection{Masses and hierarchy}

The fermionic superpartners of the off-diagonal mesons form three
Dirac pairs with masses set by the inverted seesaw:

\begin{center}
\begin{tabular}{lccc}
\toprule
Mesino pair & Spectator quark & $m_{\text{quark}}$ & Mesino mass \\
\midrule
$\psi(M^u_d, M^d_u)$ & $s$ & $93.4\MeV$ & $11\eV$ \\
$\psi(M^u_s, M^s_u)$ & $d$ & $4.67\MeV$ & $230\eV$ \\
$\psi(M^d_s, M^s_d)$ & $u$ & $2.16\MeV$ & $490\eV$ \\
\bottomrule
\end{tabular}
\end{center}

The hierarchy is inverted relative to quark masses: the mesino pair
involving the lightest quarks ($u$, $d$) has the \emph{lowest} mass
(spectator = $s$, the heaviest).  The absolute scale is
$m_\psi \sim \mu \Lambda^2 / (m_i m_j) \sim 10\eV$.

\subsection{Long-range force from mesino exchange}

The mesino Compton wavelengths define Yukawa force ranges:
\begin{center}
\begin{tabular}{lccc}
\toprule
Mesino & Mass & $\hbar c / m_\psi$ & Range \\
\midrule
$\psi(ud)$ & $11\eV$ & $1.8 \times 10^7\,\text{fm}$ & $1.8\,\text{cm}$ \\
$\psi(us)$ & $230\eV$ & $8.6 \times 10^5\,\text{fm}$ & $0.86\,\text{mm}$ \\
$\psi(ds)$ & $490\eV$ & $4.0 \times 10^5\,\text{fm}$ & $0.40\,\text{mm}$ \\
\bottomrule
\end{tabular}
\end{center}

These ranges (sub-mm to cm) are in the regime probed by torsion-balance
fifth-force experiments.  The E\"ot-Wash experiment~\cite{Kapner2007}
bounds new Yukawa forces at $\lambda = 1\,\text{mm}$ to coupling
$|\alpha| \lesssim 3 \times 10^{-4}$ relative to gravity.

The mesino--nucleon coupling is suppressed by the ISS seesaw
($\sim m_\psi / \Lambda$), but whether the resulting effective
coupling exceeds the E\"ot-Wash bound requires a detailed calculation
of the mesino--nucleon vertex through the confined phase---a task
analogous to computing the pion--nucleon form factor from QCD.

%======================================================================
\section{The Seiberg-dual sector: down-type quark masses}
\label{sec:dualkoide}
%======================================================================

The Seiberg-magnetic dual of the $\mathrm{SU}(3)$ theory has meson
VEVs $\langle M^j_j \rangle \propto 1/m_j$.  The ``magnetic charges''
$\xi_j = 1/\sqrt{m_j}$ are the natural variables, and the Koide
ratio $Q(\xi_d, \xi_s, \xi_b) = 0.6652$ deviates from $2/3$ by only
$0.22\%$.  Setting this exactly to $2/3$ and using one mass as input,
the other two follow from the Koide quadratic:

\begin{center}
\begin{tabular}{lccc}
\toprule
Input & Predicted & PDG 2024 & Deviation \\
\midrule
$(m_s, m_b)$ & $m_d = 4.60\MeV$ & $4.67 \pm 0.48\MeV$ & $0.1\sigma$ \\
$(m_d, m_b)$ & $m_s = 95.0\MeV$ & $93.4 \pm 0.8\MeV$ & $1.9\sigma$ \\
$(m_d, m_s)$ & $m_b = 4562\MeV$ & $4183 \pm 30\MeV$ & $12.6\sigma$ \\
\bottomrule
\end{tabular}
\end{center}

The $m_b$ prediction fails because the Koide inversion amplifies the
$0.22\%$ $Q$-deviation into a large absolute error at extreme mass
hierarchies.  The $m_d$ and $m_s$ predictions are within $2\sigma$,
consistent with the magnetic interpretation: the down-type quarks
organize as a near-Koide triple in the meson-condensate variables
$1/\sqrt{m_j}$.

%======================================================================
\section{Summary of predictions and experimental programme}
\label{sec:summary}
%======================================================================

\begin{table}[ht]
\centering
\small
\begin{tabular}{lllll}
\toprule
Quantity & sBootstrap & Experiment & Discrepancy & Section \\
\midrule
$\fpi^{\text{SUSY}}$ & $91.19\MeV$ & $92.07 \pm 0.57\MeV$ & $1.0\%$
  & \ref{sec:scale} \\
$m_\pi$ (off-diag.) & $\sqrt{2}\,\fpi = 130\MeV$ & $135$--$140\MeV$ & $4$--$7\%$
  & \ref{sec:mesons} \\
$|\qqbar|^{1/3}$ & $276\MeV$ & $\sim 260\MeV$ & $6\%$
  & \ref{sec:mesons} \\
GT discrepancy & $0.35\%$ & --- & $3.8\times$ improved
  & \ref{sec:nuclear} \\
$g_{\pi NN}$ (from GT) & $13.12$ & $13.17 \pm 0.06$ & $0.3\%$
  & \ref{sec:nuclear} \\
$m_\mu$ & $104.9\MeV$ & $105.66\MeV$ & $0.7\%$
  & \ref{sec:muon} \\
$m_\sigma^{\text{sB}}$ & $92\MeV$ & $f_0(500) \sim 450\MeV$ & open
  & \ref{sec:mesons} \\
$\text{STr}[M^2]$ & $18\,\fpi^2$ & not yet measured & prediction
  & \ref{sec:supertrace} \\
$m_d$ (dual Koide) & $4.60\MeV$ & $4.67 \pm 0.48\MeV$ & $0.1\sigma$
  & \ref{sec:dualkoide} \\
$m_s$ (dual Koide) & $95.0\MeV$ & $93.4 \pm 0.8\MeV$ & $1.9\sigma$
  & \ref{sec:dualkoide} \\
\bottomrule
\end{tabular}
\caption{Summary of nuclear and hadronic predictions.}
\end{table}

The sBootstrap Lagrangian produces a coherent set of predictions at
the nuclear scale:
\begin{enumerate}
\item The Goldberger--Treiman relation, with $\fpi^{\text{SUSY}}$
  instead of $\fpi^{\text{PDG}}$, is $3.8\times$ more accurate.
  This is the cleanest nuclear test.

\item The pion--nucleon coupling $g_{\pi NN} = 13.12$ is a parameter-free
  output of the GT relation combined with the SUSY identification
  $\fpi^2 = m_\pi^2 - m_\mu^2$.  It matches experiment to $0.3\%$.

\item The chiral condensate $|\qqbar|^{1/3} = 276\MeV$ is a direct
  consequence of the GOR relation with SUSY inputs.

\item The muon mass is not free: $m_\mu = \sqrt{m_\pi^2 - \fpi^2}
  = 104.9\MeV$ ($0.7\%$ accuracy), connecting leptonic and hadronic
  scales through SUSY breaking.

\item The diagonal scalar mode at $92\MeV$ provides a medium-range
  nuclear attraction at $2.2\,\text{fm}$---a testable prediction for
  the $NN$ potential.

\item The mesino spectrum ($11$--$490\eV$) predicts a new Yukawa force
  at sub-mm to cm range, accessible to fifth-force experiments.

\item The pion decay phase space is $\propto \fpi^4$ in the SUSY limit:
  the pion lifetime is controlled by the fourth power of the
  SUSY-breaking scale.

\item The $dt$ and $pp$ astrophysical $S$-factors are in principle
  computable from the sBootstrap potential without fitted nuclear
  parameters.  The longer-range sigma at $92\MeV$ modifies the
  wave function overlap at the nuclear surface.

\item The Hoyle state near-degeneracy ($\Delta E = 379\keV$ above
  $3\alpha$ threshold) may be a consequence rather than a fine-tuning:
  the sBootstrap sigma range ($2.16\,\text{fm}$) matches the
  $\alpha$-cluster separation, and the Yukawa factor at that distance
  is $1{,}100\times$ larger than the standard sigma.

\item The nuclear equation of state at ICF-relevant compressions
  ($\rho \sim 300$--$1000\,\text{g/cm}^3$, internuclear distance
  $\sim 1.8\,\text{fm}$) enters the sBootstrap sigma exchange range
  and is expected to be softer than standard Skyrme parameterizations.
\end{enumerate}

The most urgent experimental tests are: (a)~measuring or bounding
a scalar-isoscalar state below $100\MeV$; (b)~improving fifth-force
limits at $0.1$--$10\,\text{mm}$ range; (c)~lattice QCD computation
of MSbar quark masses to $1\%$ to test the dual Koide; (d)~a
precision comparison of $g_{\pi NN}$ and $g_A$ to test whether the
``correct'' $\fpi$ for the GT relation is $\sqrt{m_\pi^2 - m_\mu^2}$;
(e)~an R-matrix calculation of $S(E)$ for $d(t,n)\alpha$ with the
sBootstrap potential to test the low-energy extrapolation;
(f)~an $\alpha$-cluster model calculation of the Hoyle state energy
with the $92\MeV$ sigma exchange.

\begin{thebibliography}{99}

\bibitem{Rivero2005}
A.~Rivero, ``An interpretation of scalars in SO(32),''
Eur.\ Phys.\ J.\ C \textbf{84}, 1058 (2024),
arXiv:2407.05397 [hep-ph].

\bibitem{paperII}
A.~Rivero, ``The Koide formula and quark-lepton mass relations:
evidence and statistical assessment,'' in preparation.

\bibitem{paperIII}
A.~Rivero, ``Three families from SUSY confinement: spectrum and mass
predictions of an $\mathcal{N}=1$ $\mathrm{SU}(3)\times\mathrm{SU}(2)$
theory,'' in preparation.

\bibitem{Miyazawa}
H.~Miyazawa, ``Baryon Number Changing Currents,''
Prog.\ Theor.\ Phys.\ \textbf{36}, 1266 (1966).

\bibitem{CattoGursey}
S.~Catto and F.~G\"ursey, ``Algebraic treatment of effective
supersymmetry,'' Nuovo Cim.\ A \textbf{86}, 201 (1985).

\bibitem{Seiberg1995}
N.~Seiberg, ``Electric--magnetic duality in supersymmetric
non-abelian gauge theories,'' Nucl.\ Phys.\ B \textbf{435}, 129
(1995).

\bibitem{ISS2006}
K.~Intriligator, N.~Seiberg, D.~Shih, ``Dynamical SUSY breaking
in meta-stable vacua,'' JHEP \textbf{0604}, 021 (2006),
arXiv:hep-ph/0602239.

\bibitem{ORaifeartaigh}
L.~O'Raifeartaigh, ``Spontaneous symmetry breaking for chiral
scalar superfields,'' Nucl.\ Phys.\ B \textbf{96}, 331 (1975).

\bibitem{GOR1968}
M.~Gell-Mann, R.~J.~Oakes, B.~Renner, ``Behavior of current
divergences under $SU(3) \times SU(3)$,''
Phys.\ Rev.\ \textbf{175}, 2195 (1968).

\bibitem{GT1958}
M.~L.~Goldberger and S.~B.~Treiman, ``Decay of the pi meson,''
Phys.\ Rev.\ \textbf{110}, 1178 (1958).

\bibitem{DGMLY1967}
T.~Das, G.~S.~Guralnik, V.~S.~Mathur, F.~E.~Low, J.~E.~Young,
``Electromagnetic mass difference of pions,''
Phys.\ Rev.\ Lett.\ \textbf{18}, 759 (1967).

\bibitem{PDG2024}
R.~L.~Workman et~al.\ (Particle Data Group), ``Review of Particle
Physics,'' Prog.\ Theor.\ Exp.\ Phys.\ \textbf{2022}, 083C01 (2022);
2024 update.

\bibitem{Vesman1967}
E.~A.~Vesman, ``Concerning one possible mechanism of production of
the mesic-molecular ion $dd\mu^+$,''
JETP Lett.\ \textbf{5}, 91 (1967).

\bibitem{JonesMCF1986}
S.~E.~Jones, ``Muon-catalysed fusion revisited,''
Nature \textbf{321}, 127 (1986).

\bibitem{Kapner2007}
D.~J.~Kapner et~al., ``Tests of the gravitational inverse-square law
below the dark-energy length scale,''
Phys.\ Rev.\ Lett.\ \textbf{98}, 021101 (2007).

\bibitem{Hoyle1954}
F.~Hoyle, ``On nuclear reactions occurring in very hot stars.
I. The synthesis of elements from carbon to nickel,''
Astrophys.\ J.\ Suppl.\ \textbf{1}, 121 (1954).

\bibitem{OberhumerCSoto2000}
H.~Oberhummer, A.~Cs\'ot\'o, H.~Schlattl,
``Stellar production rates of carbon and its abundance in the
universe,'' Science \textbf{289}, 88 (2000),
arXiv:astro-ph/0007178.

\bibitem{AdelbergerSFactor2011}
E.~G.~Adelberger et~al., ``Solar fusion cross sections.
II. The $pp$ chain and CNO cycles,''
Rev.\ Mod.\ Phys.\ \textbf{83}, 195 (2011),
arXiv:1004.2318.

\end{thebibliography}

\tableofcontents

\end{document}
